% ============================================================================
% Fichier  : partie5/P05_C24_communiquer.tex
% Titre    : Communiquer en Crise
% Auteur   : MINKA MI NGUIDJOÏ Thierry Emmanuel
% Version  : 1.0 -- novembre 2026
% ============================================================================
% Sources : Coombs (2007), théorie situationnelle de la communication de
% crise ; ISO 22361:2022. Opérations GAROUA et MOUNGO : cas fictifs.
% ============================================================================

\chapter{Communiquer en Crise}
\label{chap:communiquer}
\chaptermark{Communiquer en crise}

\epigraph{%
    « Dites vite ce que vous savez, dites ce que vous ne savez pas,\\
    et dites quand vous en direz plus. »%
}{Principe du cours}

\niveauChapitre{Intermédiaire à expert}{Chapitres~\ref{chap:perception},
\ref{chap:cellule} et~\ref{chap:decider}.}

\begin{boxobjectifs}
\begin{itemize}
    \item Appliquer les principes de la communication de crise.
    \item Identifier les parties prenantes et l'ordre dans lequel les
          informer.
    \item Choisir une stratégie de réponse selon la théorie de Coombs.
    \item Rédiger un message d'attente, un communiqué et une fiche de
          questions et réponses.
    \item Gérer les réseaux sociaux, le porte-parole et les personnes dont
          les données sont touchées.
\end{itemize}
\end{boxobjectifs}

\bigskip

% ============================================================================
\section*{Question directrice : pourquoi la photographie du 6 mars a-t-elle parlé avant LiZbiZ~?}
% ============================================================================

À 18~h~30 le 6~mars, un salarié publie la photographie du message de rançon.
LiZbiZ ne dit rien avant le lendemain 9~h. Entre-temps, les clients
apprennent l'attaque par les réseaux sociaux, et les journalistes appellent
un standard sans consigne. Ce chapitre montre comment parler tôt, juste et
d'une seule voix.

\medskip
\textbf{Plan du chapitre.}
\begin{enumerate}
    \item Les principes (section~\ref{sec:c24-principes}).
    \item Les parties prenantes (section~\ref{sec:c24-parties}).
    \item La théorie de Coombs (section~\ref{sec:c24-coombs}).
    \item Les messages (section~\ref{sec:c24-messages}).
    \item Porte-parole et réseaux sociaux (section~\ref{sec:c24-canaux}).
\end{enumerate}

\separateur

% ============================================================================
\section{Les Principes}
\label{sec:c24-principes}
% ============================================================================

\begin{tableaucours}{L{3.0cm}L{9.2cm}}{%
    \tableauHeader{Principe} &
    \tableauHeader{Application}
}
Tôt & Un premier message dans l'heure qui suit la décision de crise, même
bref \\
Exact & Ne dire que ce qui est établi~; dire ce que l'on ne sait pas \\
Régulier & Annoncer quand viendra la prochaine information, et tenir
l'heure \\
Cohérent & Une seule voix, un seul contenu pour tous les canaux \\
Empathique & Parler d'abord des personnes touchées, ensuite de
l'entreprise \\
Interne d'abord & Les salariés apprennent la crise par l'entreprise, pas
par la presse \\
\end{tableaucours}
\vspace{-0.8em}
\begin{center}{\small\itshape Tableau~24.1. Les principes de la communication de crise}\end{center}

\begin{boxalert}
\textbf{Violation typique.} Se taire « tant qu'on ne sait pas tout ».

\smallskip
\textbf{Règle~:} le silence est rempli par d'autres. On publie un
\emph{message d'attente}~: ce qui s'est passé, ce que l'on fait, quand on
en dira plus. On n'attend pas la cause pour parler de la situation.
\end{boxalert}

% ============================================================================
\section{Les Parties Prenantes}
\label{sec:c24-parties}
% ============================================================================

\begin{tableaucours}{L{3.4cm}L{5.0cm}L{3.8cm}}{%
    \tableauHeader{Partie prenante} &
    \tableauHeader{Ce qu'elle attend} &
    \tableauHeader{Canal}
}
Salariés & Que faire, que dire, leur emploi est-il menacé & Réunion,
message interne \\
Clients majeurs & Leurs expéditions, leurs données, les délais & Appel du
directeur commercial \\
Clients touchés par une fuite & Quelles données, quels risques, quoi faire
& Lettre individuelle \\
Autorités & Faits, mesures, respect des obligations & Courrier, selon le
conseil juridique \\
Partenaires et prestataire & Coordination, obligations du contrat &
Cellule de crise \\
Médias & Faits, responsabilité, mesures & Communiqué, porte-parole \\
\end{tableaucours}
\vspace{-0.8em}
\begin{center}{\small\itshape Tableau~24.2. Les parties prenantes d'une crise chez LiZbiZ}\end{center}

% ============================================================================
\section{La Théorie de Coombs}
\label{sec:c24-coombs}
% ============================================================================

Timothy Coombs propose une théorie situationnelle de la communication de
crise~: la réponse doit être adaptée à la \emph{responsabilité} que les
parties prenantes attribuent à l'organisation~\cite{coombs2007}.

\begin{tableaucours}{L{3.0cm}L{4.6cm}L{4.6cm}}{%
    \tableauHeader{Groupe de crise} &
    \tableauHeader{Exemples} &
    \tableauHeader{Responsabilité attribuée}
}
Victime & Catastrophe naturelle, rumeur, malveillance extérieure & Faible \\
Accident & Défaillance technique, accident industriel & Moyenne \\
Évitable & Erreur humaine, négligence, faute de la direction & Forte \\
\end{tableaucours}
\vspace{-0.8em}
\begin{center}{\small\itshape Tableau~24.3. Les groupes de crise (d'après Coombs, 2007)}\end{center}

Les stratégies de réponse vont de \emph{nier} (contester la crise ou la
responsabilité) à \emph{diminuer} (justifier, minimiser), \emph{reconstruire}
(s'excuser, réparer, compenser), complétées par le \emph{renforcement}
(rappeler les mérites passés). Plus la responsabilité attribuée est forte,
plus il faut reconstruire. Un passé de crises semblables et une mauvaise
réputation antérieure aggravent l'attribution.

\begin{boxscenario}{Opération GAROUA. Victime ou négligente~?}
\textbf{Constatation.} LiZbiZ est victime d'une attaque criminelle, mais
la presse apprend le 21~mars que l'accès du prestataire n'avait pas de
second facteur, et qu'un audit l'avait recommandé en 2023.

\medskip
\textbf{Analyse} (confiance moyenne). La crise passe du groupe « victime »
au groupe « évitable ». La stratégie « nous sommes victimes » ne suffit
plus~: il faut reconnaître le manquement, dire ce qui a changé et
proposer une réparation aux clients touchés. Insister sur la seule qualité
de victime aurait été perçu comme une dénégation.
\end{boxscenario}

% ============================================================================
\section{Les Messages}
\label{sec:c24-messages}
% ============================================================================

\begin{boxprocedure}
\textbf{Le message d'attente}
\begin{enumerate}
    \item Ce qui s'est passé, en une phrase, sans spéculer.
    \item Ce que l'entreprise fait pour les personnes touchées.
    \item Ce que l'on ne sait pas encore.
    \item Ce que l'on demande (ne pas répondre à des messages suspects,
          appeler un numéro dédié).
    \item Quand viendra la prochaine information.
\end{enumerate}
\end{boxprocedure}

\textbf{Exemple.} « LiZbiZ subit depuis hier soir une attaque informatique
qui perturbe le suivi des expéditions à Douala et à Garoua. Les livraisons
continuent~; nos équipes informent chaque client concerné. Nous ne savons
pas encore si des données ont été dérobées. Pour toute question, appelez
le numéro dédié. Nous ferons un nouveau point demain à 12~h. »

La \termecle{fiche de questions et réponses} prépare, pour le porte-parole
et le standard, les réponses aux questions prévisibles (« avez-vous payé~?
des données ont-elles été volées~? quand la situation sera-t-elle
rétablie~? »), et la réponse à donner quand on ne sait pas.

\begin{boxscenario}{Opération MOUNGO. Les clients des colis endommagés}
\textbf{Constatation.} 15~août, 11~h~: lettre aux clients dont des colis
ont été endommagés (41 sur 47~; les six autres le 17~août, faute d'adresse
à jour)~; elle indique les colis concernés, la procédure d'indemnisation
et un interlocuteur nommé. Aucune publication de presse.

\medskip
\textbf{Analyse} (confiance élevée). Crise du groupe « victime » (pluies
exceptionnelles), traitée par une communication ciblée, rapide, orientée
réparation~: aucune escalade médiatique n'a suivi.
\end{boxscenario}

% ============================================================================
\section{Porte-Parole et Réseaux Sociaux}
\label{sec:c24-canaux}
% ============================================================================

\begin{tableaucours}{L{3.4cm}L{8.8cm}}{%
    \tableauHeader{Sujet} &
    \tableauHeader{Règle}
}
Porte-parole & Un seul, désigné, formé~; le DG pour les annonces majeures,
la communication pour le quotidien \\
Salariés & Consigne écrite~: ne pas publier, renvoyer les questions vers la
communication~; expliquer pourquoi \\
Réseaux sociaux & Surveiller~; répondre aux questions sur le même contenu
que le communiqué~; corriger les erreurs factuelles sans polémique \\
Données personnelles & Message individuel aux personnes concernées, selon
l'avis juridique (Chapitre~\ref{chap:droit}) \\
Ce qu'on ne dit pas & Les détails techniques utiles à l'attaquant, les
noms de salariés, les spéculations \\
\end{tableaucours}
\vspace{-0.8em}
\begin{center}{\small\itshape Tableau~24.4. Porte-parole et canaux}\end{center}

\begin{boxmetier}{Chargé de communication de crise}
\textbf{Préparer les modèles à froid.} Mme~ABOMO tient prêts des modèles de
message d'attente, de communiqué, de lettre aux clients et de fiche de
questions et réponses pour les cinq scénarios majeurs du registre~; en
crise, il ne reste qu'à les compléter et à les faire valider.
\end{boxmetier}

\begin{boxobjet}{R-21 · Atteinte à la réputation après une crise}
\textbf{Scénario~:} du fait d'une communication tardive ou défensive, les
clients et le public attribuent à LiZbiZ une forte responsabilité, ce qui
entraîne des pertes de clients. \textbf{Cotation~:} V=3, I=3.
\textbf{Traitement~:} modèles préparés, porte-parole formé, message
d'attente préapprouvé. \textbf{Propriétaire~:} Mme~ABOMO.
\textbf{Indicateur~:} délai entre la décision de crise et le premier message
lors des exercices.
\end{boxobjet}

\separateur

\begin{boxresume}
\textbf{Ce qu'il faut retenir de ce chapitre~:}
\begin{itemize}
    \item \textbf{Tôt, exact, régulier, cohérent, empathique, interne
          d'abord.}
    \item \textbf{Chaque partie prenante} a ses attentes et son canal.
    \item \textbf{Coombs}~: adapter la réponse à la responsabilité
          attribuée~; reconstruire quand elle est forte.
    \item \textbf{Le message d'attente} dit ce qui s'est passé, ce que l'on
          fait, ce que l'on ignore, quand on en dira plus.
    \item \textbf{Un porte-parole}, des consignes aux salariés, des modèles
          prêts.
\end{itemize}
\end{boxresume}

\begin{boxinfo}
\textbf{Points de contrôle de l'élève.} Je sais~: énoncer les principes~;
ordonner les parties prenantes~; situer une crise dans les groupes de
Coombs~; rédiger un message d'attente et une fiche de questions et
réponses~; fixer les règles du porte-parole.
\end{boxinfo}

\bigskip

\begin{boxexo}[title={\ding{45}\quad Exercice~24.1. Quel groupe~? \niveaubase}]
Situez dans les groupes de Coombs~: l'inondation de Bonabéri~; le
rançongiciel avant le 21~mars~; le rançongiciel après le 21~mars.
\end{boxexo}

\begin{boxexo}[title={\ding{45}\quad Exercice~24.2. Les principes \niveaubase}]
Quels principes du Tableau~24.1 LiZbiZ a-t-elle méconnus le 6~mars~?
\end{boxexo}

\begin{boxexo}[title={\ding{45}\quad Exercice~24.3. L'ordre \niveaubase}]
Dans quel ordre informer, le 6~mars~: les salariés, les clients majeurs,
la presse, le conseil d'administration~? Justifiez.
\end{boxexo}

\begin{boxexo}[title={\ding{45}\quad Exercice~24.4. Un message d'attente \niveaumoyen}]
Rédigez le message d'attente aux salariés de LiZbiZ, le 6~mars à 9~h.
\end{boxexo}

\begin{boxexo}[title={\ding{45}\quad Exercice~24.5. Questions et réponses \niveaumoyen}]
Rédigez la fiche de questions et réponses du standard pour le 7~mars~:
huit questions et leurs réponses.
\end{boxexo}

\begin{boxexo}[title={\ding{45}\quad Exercice~24.6. Le 21 mars \niveaumoyen}]
Rédigez le communiqué du 21~mars, après les révélations de la presse, selon
une stratégie de reconstruction.
\end{boxexo}

\begin{boxexo}[title={\ding{45}\quad Exercice~24.7. La lettre \niveaumoyen}]
Rédigez la lettre aux clients dont les données ont été publiées le
19~mars, en tenant compte des questions à soumettre au conseil juridique.
\end{boxexo}

\begin{boxexo}[title={\ding{45}\quad Exercice~24.8. Ce qu'on ne dit pas \niveaumoyen}]
Un journaliste demande « par quel serveur l'attaquant est-il entré~? ».
Que répond le porte-parole, et pourquoi~?
\end{boxexo}

\begin{boxexo}[title={\ding{45}\quad Exercice~24.9. Un plan de communication \niveauexpert}]
Rédigez le plan de communication de crise de LiZbiZ~: rôles, validation,
parties prenantes, canaux, modèles, consignes aux salariés.
\end{boxexo}

\begin{boxexo}[title={\ding{45}\quad Exercice~24.10. Une banque \niveauexpert}]
Le système central de la BANQUE SANAGA est arrêté depuis quatre heures et
les distributeurs ne fonctionnent plus. Rédigez le message d'attente aux
clients et précisez par quels canaux le diffuser.
\end{boxexo}

\bigskip

\begin{boxinfo}
\textbf{Pour aller plus loin.}
\begin{itemize}
    \item W.~T.~Coombs, « Protecting Organization Reputations During a
          Crisis »~\cite{coombs2007}.
    \item ISO~22361~\cite{iso22361_2022}.
\end{itemize}
\end{boxinfo}

\begin{boxinfo}
\textbf{Transition.} Cellule, décision et communication s'entraînent. Le
Chapitre~\ref{chap:exercices-crise} présente les exercices de crise.
\end{boxinfo}
